\properlane{promise-prayer-and-pledge}{The promised salvation is received through prayer and mercy}
\subsection{Why the promised return still calls for prayer}
God's promise in Jeremias is followed by the people's calling upon him,
praying and seeking with the whole heart. The Introit keeps that sequence
in abbreviated form. It also promises regathering from the places of
captivity. The question is why a settled divine promise should be
accompanied by a command to ask for its fulfilment.

Jerome answers at the point where he comments on the promise of hearing.
God could do what he promised without the captives' invocation. He
nevertheless exhorts them to prayer so that they may receive the promised
good. Jerome's reasoning distinguishes the giver's power from the
recipient's obedience. God's sufficiency gives no reason to abandon
prayer when God himself has commanded it.\footnote{Jerome,
\work{Commentary on Jeremias}, V, on Jer 29:10--14, PL 24,
cols.~893--894.}

The historical situation makes this obedience demanding. The exiles
are offered an attractive rival account of their future: return will
come soon. Jerome opposes that assurance to God's stated time and to
his command to build, plant and raise families. To seek God with the
whole heart includes accepting the life that his word requires during
exile. The promise can produce faithful petition because God has
spoken, even when present circumstances have not yet changed.

Jerome also explains the return anagogically, in relation to the end
of the world's confusion and the rest God grants his people. Anagogy
concerns final fulfilment. Here it extends the expectation of return
toward a consummation beyond present labour. His explanation begins
with the historical exiles and their promised restoration. It then
considers how God's people can await the completion of a promised good
while living through trial.

Psalm 84 beside Jeremias gives another form of the difficulty. God
has restored captivity, yet the psalm asks him to turn and revive the
people. Augustine reads the completed language prophetically and
explains the captivity through sin and pardon. He asks why the same
speaker both declares the turning accomplished and prays for it.
The prophet sees the promised act and asks that it come.\footnote{Augustine,
\work{Exposition of Psalm 84}, §§2--6.} Thanksgiving and petition can
therefore concern the same grace from different positions: confidence
in the promise and need for its fulfilment.

This helps explain the Collect's continuing protection. God's family
has grounds for confidence precisely in heavenly grace, on which it
says it solely relies. Dependence persists while the gift is enjoyed.
A person can thank God for preserving him and ask to be preserved
again. In this interpretation, the relation to Jeremias concerns the
recipient's continued turning to the giver. It does not make the
Collect a promise about a particular exile or a date of return.

\subsection{The sinner who can already cry}
Psalm 129 asks a still sharper question about the petitioner. The
voice comes from the depths. Augustine recognises the mortal life
and its burden of sin in that depth, but begins his exposition with
Jonas, whose prayer could not be contained by sea or fish. God is
present to the heart that prays. The cry can reach him from the
condition in which the petitioner now lives.\footnote{Augustine,
\work{Exposition of Psalm 129}, §§1--2.}

Augustine distinguishes someone who acknowledges the depth from
someone who has ceased to cry. Prosperity can conceal the latter's
danger. A person may reason that continued success after grave sins
proves that God takes no account of them. The lack of immediate
punishment then becomes an excuse to remain unchanged. The one who
cries has begun to reject that conclusion. He recognises a need for
rescue and expects God to hear.

The psalm's question about iniquities prevents confidence from being
based on a spotless record. It asks who could stand if the Lord
marked sins strictly. Augustine emphasises the universal reach of
that question: the speaker does not merely say that he himself could
not stand. The cry is possible because propitiation is with God.
Thus the person praying acknowledges guilt and hopes in mercy in
the same act.

Bellarmine develops the opening by distinguishing the depth of misery
from the depth of inward attention. A prayer can come from calamity
and also be uttered from a heart seriously intent on God. The
external need does not by itself make a petition devout. The
petitioner must give his heart to the cry.\footnote{Robert Bellarmine,
\work{Commentary on the Book of Psalms}, Ps 129, explanations of
verses 1--2, John O'Sullivan's abridged translation (1866).}
This explains why Jerome's whole-hearted seeking and the psalm's
humble cry can be considered together: both describe the recipient's
real turning toward God.

The Alleluia's acclamation does not cancel its penitential petition.
Hope is already present in the asking. At the Offertory the repeated
cry leads into the Secret's explicit request for pardon and guidance.
Read together, these prayers describe offerers whose confidence
includes acknowledgement of their need. They ask for their hearts
to be directed because mercy concerns what they will become as well
as what they have done.

\subsection{Propitiation and the risen Head}
Augustine asks what the psalm's propitiation is and answers by naming
the sacrifice offered for sinners. Innocent blood redeems the guilty
captives. The sinner's hope rests in Christ's act, rather than in an
ability to erase his own record. Augustine then describes God's merciful
law as pardoning the past and warning concerning the future. Forgiveness
and the instruction of the forgiven belong together.\footnote{Augustine,
\work{Exposition of Psalm 129}, §3.}

This account gives a precise Christological relation to the Secret.
The actual prayer calls the offerings a sacrifice of propitiation
and asks absolution. Augustine explains the propitiation of the
psalm by Christ's sacrifice; the prayer addresses the offerers' sins.
One can therefore understand their request as confidence in the mercy
Christ has made available. The further request for direction corresponds
to a forgiven life that still needs warning and guidance.

Augustine follows the psalm's waiting to the morning watch and Christ's
resurrection. His argument depends on the relation of Head and body.
The flesh Christ took from humanity has risen; what has occurred in
the Head gives the members hope for their own resurrection. They
have received forgiveness, yet they have not already undergone the
bodily resurrection for which they hope.\footnote{Augustine,
\work{Exposition of Psalm 129}, §§4--8.}

Augustine asks why Christ's resurrection should give another human
being reason to expect the same. His answer is the flesh Christ
assumed and offered. The priest received from humanity what he would
offer for it. Its resurrection as first-fruits gives the body a real
ground of hope. Sin does not defeat that hope, because the psalm
promises Israel redemption from all iniquities.

The Postcommunion explicitly distinguishes a pledge received from
salvation's effect still sought. That distinction permits a comparison
with Augustine's account of forgiven members awaiting the completion
already manifested in their Head. The gift received through the
mysteries can awaken a stronger desire for its completion. Petition
after reception expresses the saving end for which reception is given.
It also guards against treating a present gift as exemption from
perseverance.

The Gospel supplies another description of the interval: the crop
remains until harvest. The owner has promised an end without authorising
the servants to perform it now. Within this interpretation, the point
of comparison is the expectation of a divine act whose time and
execution belong to God. The actual harvest includes both gathering
and judgement. The salvation for which the Church prays is a serious
good to be sought, with repentance, before that end.

\subsection{What Christ's assurance permits the petitioner to expect}
The Communion's wording about whatever is asked can seem to make
confidence the only condition. Mark's surrounding instruction names
faith in God and forgiveness of an offender. Bede develops that
context by explaining faith as working through charity. Those who
have not attained its perfection can still ask forgiveness and eternal
life, provided they forgive others. His account gives the uncertain
petitioner an actual prayer to make and a relation to the neighbour
to mend.\footnote{Bede, \work{Exposition of Mark}, on Mark 11:22--26,
paragraphs \textit{Et respondens Jesus} and \textit{Et cum stabitis ad orandum}.}

Aquinas distinguishes two effects of prayer. Its merit arises from
charity; its obtaining what is asked, often called impetration,
arises from God's gracious gift. Faith approaches God's power and
mercy, humility acknowledges need, and charity desires the good.
The distinction matters when a prayer is devout but the particular
temporal thing requested would not conduce to the person's salvation.
The act can be good without that desired outcome being granted.\footnote{Aquinas,
\work{Summa theologiae} II--II, q.~83, a.~15, corpus and replies 1--3;
English Dominican translation. This is doctrinal exposition of prayer,
not Aquinas's direct commentary on Mark 11.}

His four conditions concern asking for oneself, for things necessary
to salvation, piously and perseveringly. He also allows a delay in a
useful gift, to be granted at a fitting time. The conditions specify
the saving object and the manner of asking. They do not provide a
method of diagnosing each unanswered temporal petition as deficient
faith. Aquinas expressly recognises that a person may ask faithfully
for a present good and be mercifully left without it because God
knows what is conducive to beatitude.

He then asks whether sinners obtain anything from God. He distinguishes
the nature God loves from the sin he hates. A sinful desire is different
from a good request made by someone who has sinned. God can hear the
latter out of mercy.\footnote{Aquinas, \work{Summa theologiae}
II--II, q.~83, a.~16, corpus.} This distinction explains why the cry
from the depths need not wait until its speaker has become
self-sufficient. Asking for the pardon and life that God wills can
itself belong to the sinner's return.

\subsection{Instruction, praise and thanksgiving while fulfilment is awaited}
The Epistle makes thanksgiving habitual. The forgiven are told to
forgive; the one body is told to let Christ's word dwell richly.
The person who awaits salvation can learn, teach, admonish and sing
now. These activities address the heart that asks. They give it
words and desires formed by Christ rather than by the grievance or
fear of a particular moment.

Chrysostom's account of charity supplies the moral condition of this
common life. Aquinas's exposition of Colossians 3:16--17 explains
instruction and heartfelt song, then the direction of ordinary acts
to Christ. Its habitual intention extends thanksgiving beyond the
explicit moments of prayer.\footnote{Chrysostom, \work{Homilies on
Colossians}, VIII--IX; Aquinas, \work{Commentary on Colossians},
chapter 3, lectio 3, on 3:16--17, Dessain II, p.~442.}
One who asks for salvation's effect can already act in the name of
the Lord and give thanks through him.

The Gradual's confident praise, within Psalm 43's lament, permits
thanksgiving during suffering. Augustine hears both the martyrs'
present trial and the certainty of future deliverance. The petitioner
has reason to ask even when outward circumstances seem unchanged.
The Preface gives the corresponding actual act of thanksgiving,
always and everywhere, to the Father with the Son and Holy Spirit.
The saving gift is received from God and answered with adoration.

Jerome's exiles pray because God has promised. Augustine's sinner
cries because mercy is with the Lord. In both arguments, asking
acknowledges something already given: God's word and the possibility
of pardon. The Postcommunion goes further by naming the pledge already
received through the mysteries. The final petition seeks the completion
of a salvation whose giver has already acted.

\subsection{Four senses}
\begin{fourSenses}
\item[Literal.] God promises restoration to exiles and commands them to
call and seek. The psalmist cries, acknowledges sin and waits. Christ
commands believing and forgiving prayer; the orations ask protection,
pardon, direction and salvation's effect.
\item[Allegorical.] Christ's sacrifice is the propitiation through which
sinners approach God. His resurrection in the flesh gives the waiting
members hope in their risen Head. The mysteries' pledge belongs to
this salvation received through him.
\item[Moral.] The petitioner asks for saving good piously and perseveringly,
forgives an actual offender, accepts instruction and gives thanks in
ordinary work. Confidence acknowledges mercy and seeks the conversion
of the person who asks.
\item[Anagogical.] The promised return reaches complete redemption and
bodily resurrection. The received pledge anticipates salvation's effect;
the harvest ends the interval of waiting and fulfils Christ's judgement.
\end{fourSenses}
